\documentclass[11pt]{article}
% Shared preamble for the printable mini lessons (1.4 onward).
%   \documentclass[11pt]{article}
%   % Shared preamble for the printable mini lessons (1.4 onward).
%   \documentclass[11pt]{article}
%   % Shared preamble for the printable mini lessons (1.4 onward).
%   \documentclass[11pt]{article}
%   \input{../preamble}
% Build the key with:  pdflatex -jobname=<name>_KEY "\def\iskey{}\input{<name>}"
% Figures shared with the web page live in figures/Unit_X/<name>.tex and are
% pulled in with \fig{Unit_X/<name>}; the standalone package drops their own
% preamble, so their macros come from figures/preamble.tex, loaded here.
\input{../../figures/preamble}
\usepackage[margin=0.75in]{geometry}
\usepackage{array}
\usepackage{enumitem}
\usepackage{amsmath}
\usepackage[helvet]{sfmath}  % math in the same sans font as the text
\usepackage{standalone}
\pagestyle{empty}
\setlength{\parindent}{0pt}
\setlist[enumerate,1]{leftmargin=1.5em,itemsep=1em}
\setlist[enumerate,2]{label=(\alph*),leftmargin=2em,itemsep=2.2em}

\newif\ifkey
\ifdefined\iskey \keytrue \fi
\colorlet{keyred}{red!75!black}
\newcommand{\ans}[1]{\ifkey\leavevmode{\color{keyred}#1}\fi}
% vertical work space on the student copy, less on the key
\newcommand{\work}[1]{\ifkey\vspace{0.6em}\else\vspace{#1}\fi}
\newcommand{\fig}[1]{\input{../../figures/#1}}
% Title line: \lessontitle{Quartiles, IQR, and Box Plots}
\newcommand{\lessontitle}[1]{{\Large\bfseries #1\ifkey\ \textcolor{keyred}{KEY}\fi}\hfill Name \underline{\hspace{2in}}\par\vspace{0.6em}}
\newcommand{\partner}{\par\vspace{0.4em}Work with a partner. Compare answers when you finish each problem.\par}

% Build the key with:  pdflatex -jobname=<name>_KEY "\def\iskey{}\input{<name>}"
% Figures shared with the web page live in figures/Unit_X/<name>.tex and are
% pulled in with \fig{Unit_X/<name>}; the standalone package drops their own
% preamble, so their macros come from figures/preamble.tex, loaded here.
% Shared preamble for web figures.  Each figure is a standalone .tex that
% \input's this file; scripts/make-figures.sh turns it into an SVG.
\usepackage[T1]{fontenc}
\usepackage{helvet}
\renewcommand{\familydefault}{\sfdefault}
\usepackage{tikz}
\usetikzlibrary{decorations.pathreplacing}
% axis numbers in the text font, not math italic
\pgfkeys{/pgf/number format/.cd,assume math mode=true,1000 sep={}}

% dot plot: \dotplot{min}{max}{value/count,...}{axis label}
\newcommand{\dotplot}[4]{%
\begin{tikzpicture}[x=0.9cm,y=0.33cm]
  \draw[thick] (#1-0.5,0) -- (#2+0.5,0);
  \foreach \t in {#1,...,#2} {\draw (\t,0.25) -- (\t,-0.25) node[below] {\t};}
  \foreach \v/\c in {#3} {\ifnum\c>0 \foreach \k in {1,...,\c} {\fill (\v,\k*1.15) circle (0.16cm);}\fi}
  \node at ({(#1+#2)/2},-2.3) {#4};
\end{tikzpicture}}

% shape sketch histogram (12 bins): \shapehist{heights}{label}
\newcommand{\shapehist}[2]{%
  \foreach \h [count=\i from 0] in {#1} \draw[fill=gray!30] (\i,0) rectangle (\i+1,\h);
  \draw (-0.3,0) -- (12.3,0);
  \node[font=\bfseries] at (6,-2.6) {#2};}
\newcommand{\meanmark}[1]{\fill (#1,0) -- ++(-0.3,-0.75) -- ++(0.6,0) -- cycle;}
\newcommand{\medmark}[2]{\draw[very thick,dashed] (#1,0) -- (#1,#2);}

% number line with minor ticks: \nline{min}{max}{label step}{axis label}
\newcommand{\nline}[4]{%
  \draw[thick] (#1-0.5,0) -- (#2+0.5,0);
  \foreach \t in {#1,...,#2} \draw (\t,0.12) -- (\t,-0.12);
  \pgfmathtruncatemacro{\nxt}{#1+#3}
  \foreach \t in {#1,\nxt,...,#2} \draw (\t,0.2) -- (\t,-0.2) node[below] {\t};
  \node at ({(#1+#2)/2},-0.95) {#4};}
% box plot above the line: \bplot{min}{Q1}{med}{Q3}{max}
\newcommand{\bplot}[5]{%
  \draw[thick] (#1,1.0) -- (#2,1.0);
  \draw[thick] (#4,1.0) -- (#5,1.0);
  \draw[thick] (#1,0.75) -- (#1,1.25);
  \draw[thick] (#5,0.75) -- (#5,1.25);
  \draw[thick,fill=gray!20] (#2,0.55) rectangle (#4,1.45);
  \draw[very thick] (#3,0.55) -- (#3,1.45);}

% ---- added for 1.4 onward ----
% axis on y=0: \axis{min}{max}{label step}{tick step}{axis label}
\newcommand{\axis}[5]{%
  \draw[thick] (#1,0) -- (#2,0);
  \pgfmathsetmacro{\axT}{#1+#4}\foreach \t in {#1,\axT,...,#2} \draw (\t,0.1) -- (\t,-0.1);
  \pgfmathsetmacro{\axL}{#1+#3}\foreach \t in {#1,\axL,...,#2}
    \draw (\t,0.18) -- (\t,-0.18) node[below] {\pgfmathprintnumber{\t}};
  \node at ({(#1+#2)/2},-0.95) {#5};}
% box plot centred at height y: \bplotat{y}{min}{Q1}{med}{Q3}{max}
\newcommand{\bplotat}[6]{%
  \draw[thick] (#2,#1) -- (#3,#1);
  \draw[thick] (#5,#1) -- (#6,#1);
  \draw[thick] (#2,#1-0.22) -- (#2,#1+0.22);
  \draw[thick] (#6,#1-0.22) -- (#6,#1+0.22);
  \draw[thick,fill=gray!20] (#3,#1-0.42) rectangle (#5,#1+0.42);
  \draw[very thick] (#4,#1-0.42) -- (#4,#1+0.42);}
% an outlier dot: \outlier{x}{y}
\newcommand{\outlier}[2]{\fill (#1,#2) circle (2.4pt);}

\usepackage[margin=0.75in]{geometry}
\usepackage{array}
\usepackage{enumitem}
\usepackage{amsmath}
\usepackage[helvet]{sfmath}  % math in the same sans font as the text
\usepackage{standalone}
\pagestyle{empty}
\setlength{\parindent}{0pt}
\setlist[enumerate,1]{leftmargin=1.5em,itemsep=1em}
\setlist[enumerate,2]{label=(\alph*),leftmargin=2em,itemsep=2.2em}

\newif\ifkey
\ifdefined\iskey \keytrue \fi
\colorlet{keyred}{red!75!black}
\newcommand{\ans}[1]{\ifkey\leavevmode{\color{keyred}#1}\fi}
% vertical work space on the student copy, less on the key
\newcommand{\work}[1]{\ifkey\vspace{0.6em}\else\vspace{#1}\fi}
\newcommand{\fig}[1]{\input{../../figures/#1}}
% Title line: \lessontitle{Quartiles, IQR, and Box Plots}
\newcommand{\lessontitle}[1]{{\Large\bfseries #1\ifkey\ \textcolor{keyred}{KEY}\fi}\hfill Name \underline{\hspace{2in}}\par\vspace{0.6em}}
\newcommand{\partner}{\par\vspace{0.4em}Work with a partner. Compare answers when you finish each problem.\par}

% Build the key with:  pdflatex -jobname=<name>_KEY "\def\iskey{}\input{<name>}"
% Figures shared with the web page live in figures/Unit_X/<name>.tex and are
% pulled in with \fig{Unit_X/<name>}; the standalone package drops their own
% preamble, so their macros come from figures/preamble.tex, loaded here.
% Shared preamble for web figures.  Each figure is a standalone .tex that
% \input's this file; scripts/make-figures.sh turns it into an SVG.
\usepackage[T1]{fontenc}
\usepackage{helvet}
\renewcommand{\familydefault}{\sfdefault}
\usepackage{tikz}
\usetikzlibrary{decorations.pathreplacing}
% axis numbers in the text font, not math italic
\pgfkeys{/pgf/number format/.cd,assume math mode=true,1000 sep={}}

% dot plot: \dotplot{min}{max}{value/count,...}{axis label}
\newcommand{\dotplot}[4]{%
\begin{tikzpicture}[x=0.9cm,y=0.33cm]
  \draw[thick] (#1-0.5,0) -- (#2+0.5,0);
  \foreach \t in {#1,...,#2} {\draw (\t,0.25) -- (\t,-0.25) node[below] {\t};}
  \foreach \v/\c in {#3} {\ifnum\c>0 \foreach \k in {1,...,\c} {\fill (\v,\k*1.15) circle (0.16cm);}\fi}
  \node at ({(#1+#2)/2},-2.3) {#4};
\end{tikzpicture}}

% shape sketch histogram (12 bins): \shapehist{heights}{label}
\newcommand{\shapehist}[2]{%
  \foreach \h [count=\i from 0] in {#1} \draw[fill=gray!30] (\i,0) rectangle (\i+1,\h);
  \draw (-0.3,0) -- (12.3,0);
  \node[font=\bfseries] at (6,-2.6) {#2};}
\newcommand{\meanmark}[1]{\fill (#1,0) -- ++(-0.3,-0.75) -- ++(0.6,0) -- cycle;}
\newcommand{\medmark}[2]{\draw[very thick,dashed] (#1,0) -- (#1,#2);}

% number line with minor ticks: \nline{min}{max}{label step}{axis label}
\newcommand{\nline}[4]{%
  \draw[thick] (#1-0.5,0) -- (#2+0.5,0);
  \foreach \t in {#1,...,#2} \draw (\t,0.12) -- (\t,-0.12);
  \pgfmathtruncatemacro{\nxt}{#1+#3}
  \foreach \t in {#1,\nxt,...,#2} \draw (\t,0.2) -- (\t,-0.2) node[below] {\t};
  \node at ({(#1+#2)/2},-0.95) {#4};}
% box plot above the line: \bplot{min}{Q1}{med}{Q3}{max}
\newcommand{\bplot}[5]{%
  \draw[thick] (#1,1.0) -- (#2,1.0);
  \draw[thick] (#4,1.0) -- (#5,1.0);
  \draw[thick] (#1,0.75) -- (#1,1.25);
  \draw[thick] (#5,0.75) -- (#5,1.25);
  \draw[thick,fill=gray!20] (#2,0.55) rectangle (#4,1.45);
  \draw[very thick] (#3,0.55) -- (#3,1.45);}

% ---- added for 1.4 onward ----
% axis on y=0: \axis{min}{max}{label step}{tick step}{axis label}
\newcommand{\axis}[5]{%
  \draw[thick] (#1,0) -- (#2,0);
  \pgfmathsetmacro{\axT}{#1+#4}\foreach \t in {#1,\axT,...,#2} \draw (\t,0.1) -- (\t,-0.1);
  \pgfmathsetmacro{\axL}{#1+#3}\foreach \t in {#1,\axL,...,#2}
    \draw (\t,0.18) -- (\t,-0.18) node[below] {\pgfmathprintnumber{\t}};
  \node at ({(#1+#2)/2},-0.95) {#5};}
% box plot centred at height y: \bplotat{y}{min}{Q1}{med}{Q3}{max}
\newcommand{\bplotat}[6]{%
  \draw[thick] (#2,#1) -- (#3,#1);
  \draw[thick] (#5,#1) -- (#6,#1);
  \draw[thick] (#2,#1-0.22) -- (#2,#1+0.22);
  \draw[thick] (#6,#1-0.22) -- (#6,#1+0.22);
  \draw[thick,fill=gray!20] (#3,#1-0.42) rectangle (#5,#1+0.42);
  \draw[very thick] (#4,#1-0.42) -- (#4,#1+0.42);}
% an outlier dot: \outlier{x}{y}
\newcommand{\outlier}[2]{\fill (#1,#2) circle (2.4pt);}

\usepackage[margin=0.75in]{geometry}
\usepackage{array}
\usepackage{enumitem}
\usepackage{amsmath}
\usepackage[helvet]{sfmath}  % math in the same sans font as the text
\usepackage{standalone}
\pagestyle{empty}
\setlength{\parindent}{0pt}
\setlist[enumerate,1]{leftmargin=1.5em,itemsep=1em}
\setlist[enumerate,2]{label=(\alph*),leftmargin=2em,itemsep=2.2em}

\newif\ifkey
\ifdefined\iskey \keytrue \fi
\colorlet{keyred}{red!75!black}
\newcommand{\ans}[1]{\ifkey\leavevmode{\color{keyred}#1}\fi}
% vertical work space on the student copy, less on the key
\newcommand{\work}[1]{\ifkey\vspace{0.6em}\else\vspace{#1}\fi}
\newcommand{\fig}[1]{\input{../../figures/#1}}
% Title line: \lessontitle{Quartiles, IQR, and Box Plots}
\newcommand{\lessontitle}[1]{{\Large\bfseries #1\ifkey\ \textcolor{keyred}{KEY}\fi}\hfill Name \underline{\hspace{2in}}\par\vspace{0.6em}}
\newcommand{\partner}{\par\vspace{0.4em}Work with a partner. Compare answers when you finish each problem.\par}

% Unit summary: the unit's tools in tables, to print and keep. Same tables as
% chapters/Unit_1/Summary.qmd; edit both together. Nothing to answer, so the
% _KEY copy is the same sheet.
\newcommand{\sumhead}[1]{\par\vspace{0.7em}{\large\bfseries #1}\par\vspace{0.3em}}
\newcolumntype{L}[1]{>{\raggedright\arraybackslash}p{#1}}
\renewcommand{\arraystretch}{1.25}

\begin{document}

\lessontitle{Unit 1 Summary: Displaying One Variable}

\sumhead{Showing the Data}
\begin{tabular}{|L{1.2in}|L{2.95in}|L{2.3in}|}\hline
\textbf{Display} & \textbf{How it works} & \textbf{Example}\\ \hline
Frequency table & Lists each value and its \textbf{frequency} (how many times it happens). Make a tally mark for each value. Then count the marks. & 11 students say how many siblings they have. 4 of them have 1 sibling. Check: $2 + 4 + 3 + 1 + 1 = 11$ students.\\ \hline
Dot plot & Puts the values on a number line. Each value gets \textbf{one dot}. Dots for the same value stack up. The height of a stack is the frequency. It shows every value. & The same siblings data. The stack above 1 is the tallest, so 1 is the most common.\\ \hline
Histogram & A bar graph for number data. The values are put into groups called \textbf{bins}. Each bar is one bin. The bars touch (no gaps). It does \textbf{not} show the exact values. It is better than a dot plot when there is a lot of data, or the values are spread out. & Heights of students. The tallest bar is 170 to 180. So 7 students are from 170 cm to 180 cm tall.\\ \hline
\end{tabular}

\vspace{0.3em}
\textbf{The histogram rule:} if a value is \textbf{on a line} between two bars, it goes in the bar on the \textbf{right}. 160 goes in the 160 to 170 bar.

\vspace{0.3em}
\textbf{Every time you read a graph:} read the title, read both axis labels, then check the scale.

\sumhead{Center}
\begin{tabular}{|L{1.2in}|L{2.95in}|L{2.3in}|}\hline
\textbf{Measure} & \textbf{How to find it} & \textbf{Example: 2, 3, 8, 3, 4}\\ \hline
Mean (the average) & Add all the values. Divide by how many values there are. In symbols: $\bar{x} = \dfrac{\Sigma x}{n}$.\newline $\Sigma$ (sigma) means \textbf{add up}. $n$ is how many values there are. $\bar{x}$ (``x-bar'') is the mean. & $\Sigma x = 2 + 3 + 8 + 3 + 4 = 20$ and $n = 5$, so $\bar{x} = \dfrac{20}{5} = 4$.\\ \hline
Median (the middle) & Put the values in order. Find the middle value. Two values in the middle? The median is the number halfway between them. & In order: 2, 3, \textbf{3}, 4, 8. The median is 3. For 2, 3, 5, 8: $(3 + 5) \div 2 = 4$.\\ \hline
Mode & The value that happens \textbf{most often}. & The mode is 3.\\ \hline
\end{tabular}

\sumhead{Shape}
\begin{tabular}{|L{1.2in}|L{2.95in}|L{2.3in}|}\hline
\textbf{Shape} & \textbf{What it looks like} & \textbf{Mean and median}\\ \hline
Skewed left & Tail on the left & mean $<$ median\\ \hline
Symmetric & Both sides match & mean $\approx$ median\\ \hline
Skewed right & Tail on the right & mean $>$ median\\ \hline
\end{tabular}

\vspace{0.3em}
\textbf{The mean follows the tail.} The tail is the long, low side of a graph. The mean is pulled toward the tail. The median is not pulled much. With a long tail, the median is a better description of a \textbf{typical} value.

\newpage
\sumhead{Quartiles and Spread}
\begin{tabular}{|L{2.05in}|L{2.2in}|L{2.2in}|}\hline
\textbf{How to find it} & \textbf{Example 1: 11 values (odd)} & \textbf{Example 2: 10 values (even)}\\ \hline
1. Put the data in order. & 3, 4, 6, 7, 9, \textbf{10}, 12, 13, 15, 18, 20 & 1, 3, 4, 6, \textbf{7}, \textbf{9}, 10, 12, 14, 18\\ \hline
2. Find the \textbf{median}. It splits the data into a lower half and an upper half. & The median is 10. \textbf{Odd number of values:} leave the median out of both halves. & Two middle values, 7 and 9. The median is $(7 + 9) \div 2 = 8$.\\ \hline
3. \textbf{Q1} = the median of the lower half. & Lower half: 3, 4, \textbf{6}, 7, 9. Q1~=~6 & Lower half: 1, 3, \textbf{4}, 6, 7. Q1~=~4\\ \hline
4. \textbf{Q3} = the median of the upper half. & Upper half: 12, 13, \textbf{15}, 18, 20. Q3~=~15 & Upper half: 9, 10, \textbf{12}, 14, 18. Q3~=~12\\ \hline
Five-number summary: minimum, Q1, median, Q3, maximum & 3, 6, 10, 15, 20 & 1, 4, 8, 12, 18\\ \hline
IQR $= Q_3 - Q_1$ & $15 - 6 = 9$ & $12 - 4 = 8$\\ \hline
\end{tabular}

\vspace{0.3em}
\textbf{IQR} (interquartile range): how spread out the \textbf{middle half} of the data is. \textbf{Range} $=$ maximum $-$ minimum. One extreme value changes the range a lot. It does not move the IQR much.

\sumhead{Box Plots}
\begin{tabular}{|L{1.45in}|L{2.75in}|L{2.25in}|}\hline
\textbf{Part} & \textbf{What it shows} & \textbf{Example: test scores}\\ \hline
The box & Goes from Q1 to Q3. Its length is the IQR. It holds the middle \textbf{50\%} of the data. & Half of the students scored from 68 to 84.\\ \hline
The line inside the box & The median. About \textbf{50\%} of the data is below it. & Half scored less than 76.\\ \hline
The whiskers & Go out to the minimum and the maximum. & From 52 to 68, and from 84 to 98.\\ \hline
Each of the 4 parts & About \textbf{25\%} of the data. A long part means the values are \textbf{spread out}. It does not mean there are more of them. & The left whisker (52 to 68) is the longest part. It still holds only about 25\% of the students.\\ \hline
Above Q1 & About \textbf{75\%} of the data. & About 75\% scored more than 68.\\ \hline
\end{tabular}

\vspace{0.3em}
\textbf{Compare two box plots:} put both on the same number line. Then compare the centers and the spreads.

\vspace{0.3em}
\begin{tabular}{|L{1.25in}|L{0.95in}|L{0.95in}|L{3.05in}|}\hline
\textbf{Compare} & \textbf{Morning} & \textbf{Afternoon} & \textbf{What it means}\\ \hline
Center (median) & 72 & 78 & The afternoon class scored \textbf{higher}.\\ \hline
Spread (IQR) & $80 - 65 = 15$ & $82 - 70 = 12$ & The afternoon scores are \textbf{closer together}.\\ \hline
\end{tabular}

\sumhead{Outliers}
\begin{tabular}{|L{3.6in}|L{2.85in}|}\hline
\textbf{The 1.5 $\times$ IQR rule} & \textbf{Example: 3, 8, 9, 10, 11, 12, 13, 14, 30}\\ \hline
1. Find Q1, Q3, and the IQR. & Q1 = 8.5, Q3 = 13.5, and IQR = 5.\\ \hline
2. Lower fence $= Q_1 - 1.5 \times \text{IQR}$ & $1.5 \times 5 = 7.5$. Lower fence $= 8.5 - 7.5 = 1$\\ \hline
3. Upper fence $= Q_3 + 1.5 \times \text{IQR}$ & Upper fence $= 13.5 + 7.5 = 21$\\ \hline
4. A value below the lower fence or above the upper fence is an \textbf{outlier}. It is far away from the rest of the data. & 30 is above 21, so \textbf{30 is an outlier}. 3 is not below 1, so 3 is not an outlier.\\ \hline
\end{tabular}

\vspace{0.3em}
\textbf{On a box plot:} an outlier is a \textbf{dot}. The whisker stops at the last value that is not an outlier (here, 14).

\end{document}
